\documentclass[12pt,a4paper]{report}
\usepackage[utf8]{inputenc}
\usepackage[T1]{fontenc}
\usepackage{geometry}
\usepackage{graphicx}
\usepackage{hyperref}
\usepackage{listings}
\usepackage{xcolor}
\usepackage{booktabs}
\usepackage{longtable}
\usepackage{fancyhdr}
\usepackage{titlesec}
\usepackage{tocloft}
\usepackage{amsmath}
\usepackage{amssymb}

\geometry{margin=1in}

\definecolor{codegreen}{rgb}{0,0.6,0}
\definecolor{codegray}{rgb}{0.5,0.5,0.5}
\definecolor{codepurple}{rgb}{0.58,0,0.82}
\definecolor{backcolour}{rgb}{0.95,0.95,0.92}

\lstdefinestyle{mystyle}{
    backgroundcolor=\color{backcolour},
    commentstyle=\color{codegreen},
    keywordstyle=\color{magenta},
    numberstyle=\tiny\color{codegray},
    stringstyle=\color{codepurple},
    basicstyle=\ttfamily\footnotesize,
    breakatwhitespace=false,
    breaklines=true,
    captionpos=b,
    keepspaces=true,
    numbers=left,
    numbersep=5pt,
    showspaces=false,
    showstringspaces=false,
    showtabs=false,
    tabsize=2
}
\lstset{style=mystyle}

\hypersetup{
    colorlinks=true,
    linkcolor=blue,
    filecolor=magenta,
    urlcolor=cyan,
}

\title{
    \textbf{Conversational AI Pipeline} \\
    \Large Comprehensive Codebase Documentation \\
    \vspace{0.5cm}
    \normalsize Multi-Agent System for Autonomous Data Analysis
}
\author{Technical Documentation}
\date{\today}

\begin{document}

\maketitle
\tableofcontents
\newpage

%==============================================================================
% CHAPTER 1: EXECUTIVE SUMMARY
%==============================================================================
\chapter{Executive Summary}

\section{System Overview}
The Conversational AI Pipeline is a sophisticated multi-agent system designed for autonomous data analysis, forecasting, and machine learning tasks. The system leverages Large Language Models (LLMs) orchestrated through LangGraph to create specialized agents that collaborate in a structured pipeline.

\subsection{Key Capabilities}
\begin{itemize}
    \item \textbf{Autonomous Data Profiling}: Automatic discovery and summarization of datasets
    \item \textbf{Intelligent Task Proposal}: LLM-driven analysis to propose analytical tasks
    \item \textbf{Execution Planning}: Detailed plan generation for data processing and modeling
    \item \textbf{Method Benchmarking}: Automated comparison of multiple ML/statistical methods
    \item \textbf{Visualization Generation}: Dynamic creation of insightful visualizations
    \item \textbf{Report Generation}: Comprehensive final reports with findings
    \item \textbf{Guardrails Validation}: Statistical validation of model predictions
    \item \textbf{Conversational Interface}: Natural language interaction for users
\end{itemize}

\subsection{Technology Stack}
\begin{itemize}
    \item \textbf{Language}: Python 3.10+
    \item \textbf{LLM Framework}: LangChain, LangGraph
    \item \textbf{Data Processing}: Pandas, NumPy
    \item \textbf{Machine Learning}: Scikit-learn, XGBoost, LightGBM
    \item \textbf{Visualization}: Matplotlib, Seaborn
    \item \textbf{Web Interface}: FastAPI, WebSockets
    \item \textbf{Data Validation}: Pydantic
\end{itemize}

%==============================================================================
% CHAPTER 2: ARCHITECTURE OVERVIEW
%==============================================================================
\chapter{System Architecture}

\section{High-Level Architecture}
The system follows a multi-stage pipeline architecture where each stage is implemented as an autonomous LLM agent with specific tools and responsibilities.

\subsection{Pipeline Flow}
\begin{verbatim}
User Query → Conversation Agent → Pipeline Orchestrator
                                         ↓
    Stage 1 (Dataset Summarization)
                ↓
    Stage 2 (Task Proposal Generation)
                ↓
    Stage 3 (Execution Planning)
                ↓
    Stage 3B (Data Preparation)
                ↓
    Stage 3.5A (Method Proposal)
                ↓
    Stage 3.5B (Method Benchmarking)
                ↓
    Stage 4 (Execution)
                ↓
    Stage 5 (Visualization)
                ↓
    Stage 6 (Report Generation)
                ↓
    Stage 7 (Guardrails Validation) [Optional]
\end{verbatim}

\section{Directory Structure}
\begin{lstlisting}[language=bash]
conversational/
├── code/                    # Core agent implementations
│   ├── config.py           # Centralized configuration
│   ├── models.py           # Pydantic data models
│   ├── utils.py            # Utility functions
│   ├── master_orchestrator.py  # Pipeline coordination
│   ├── conversation_agent.py   # User interaction handler
│   ├── eda_agent.py        # Exploratory data analysis
│   ├── stage1_agent.py     # Dataset summarization
│   ├── stage2_agent.py     # Task proposal generation
│   ├── stage3_agent.py     # Execution planning
│   ├── stage3b_agent.py    # Data preparation
│   ├── stage3_5a_agent.py  # Method proposal
│   ├── stage3_5b_agent.py  # Method benchmarking
│   ├── stage4_agent.py     # Execution
│   ├── stage5_agent.py     # Visualization
│   ├── stage6_agent.py     # Report generation
│   └── stage7_guardrails_agent.py  # Validation
├── tools/                   # Stage-specific tools
│   ├── stage1_tools.py     # Dataset profiling tools
│   ├── stage2_tools.py     # Task proposal tools
│   ├── stage3_tools.py     # Planning tools
│   ├── stage3b_tools.py    # Data prep tools
│   ├── stage3_5a_tools.py  # Method proposal tools
│   ├── stage3_5b_tools.py  # Benchmarking tools
│   ├── stage4_tools.py     # Execution tools
│   ├── stage5_tools.py     # Visualization tools
│   ├── stage6_tools.py     # Report tools
│   ├── stage7_guardrails_tools.py  # Validation tools
│   ├── conversation_tools.py  # Chat interface tools
│   └── eda_tools.py        # EDA analysis tools
├── ui/                      # Web interface
│   ├── static/             # HTML, CSS, JavaScript
│   └── api.py              # Additional API endpoints
├── data/                    # Input datasets
├── output/                  # Stage outputs
├── server.py               # FastAPI server
└── run_conversational.py   # CLI entry point
\end{lstlisting}

%==============================================================================
% CHAPTER 3: CONFIGURATION
%==============================================================================
\chapter{Configuration Module}

\section{Overview}
The \texttt{config.py} module provides centralized configuration for the entire pipeline, including LLM settings, directory paths, and stage-specific parameters.

\section{LLM Configuration}
\begin{lstlisting}[language=Python]
PRIMARY_LLM_CONFIG = {
    "base_url": LLM_BASE_URL,
    "api_key": LLM_API_KEY,
    "model": "Qwen/Qwen2.5-32B-Instruct",
    "temperature": 0.0,
    "max_tokens": 4096,
}

SECONDARY_LLM_CONFIG = {
    "model": "Qwen/Qwen3-32B",
    "temperature": 0.2,
    "max_tokens": 4096,
}
\end{lstlisting}

\section{Directory Configuration}
\begin{itemize}
    \item \texttt{DATA\_DIR}: Input datasets directory
    \item \texttt{OUTPUT\_ROOT}: Base output directory
    \item \texttt{SUMMARIES\_DIR}: Stage 1 dataset summaries
    \item \texttt{STAGE2\_OUT\_DIR}: Task proposals
    \item \texttt{STAGE3\_OUT\_DIR}: Execution plans
    \item \texttt{STAGE3B\_OUT\_DIR}: Prepared data
    \item \texttt{STAGE3\_5A\_OUT\_DIR}: Method proposals
    \item \texttt{STAGE3\_5B\_OUT\_DIR}: Benchmark results
    \item \texttt{STAGE4\_OUT\_DIR}: Execution results
    \item \texttt{STAGE5\_OUT\_DIR}: Visualizations
    \item \texttt{STAGE6\_OUT\_DIR}: Final reports
    \item \texttt{STAGE7\_OUT\_DIR}: Guardrails validation
\end{itemize}

\section{Stage Iteration Limits}
\begin{lstlisting}[language=Python]
STAGE_MAX_ROUNDS = {
    "stage1": 20,    # Dataset profiling
    "stage2": 40,    # Task proposal (complex exploration)
    "stage3": 30,    # Execution planning
    "stage3b": 40,   # Data preparation
    "stage3_5a": 30, # Method proposal
    "stage3_5b": 50, # Benchmarking (multiple iterations)
    "stage4": 40,    # Execution
    "stage5": 30,    # Visualization
    "stage6": 30,    # Report generation
    "stage7": 50,    # Guardrails validation
    "eda": 40,       # EDA agent exploration
}
\end{lstlisting}

\section{DataPassingManager}
The \texttt{DataPassingManager} class provides robust data passing between pipeline stages with:
\begin{itemize}
    \item Atomic file operations using temporary files
    \item SHA256 checksum verification
    \item Automatic retry on failures
    \item Structured metadata for all artifacts
\end{itemize}

%==============================================================================
% CHAPTER 4: DATA MODELS
%==============================================================================
\chapter{Data Models}

\section{Overview}
The \texttt{models.py} module defines all Pydantic data structures used throughout the pipeline, ensuring type safety and validation.

\section{Enumerations}
\subsection{TaskCategory}
\begin{lstlisting}[language=Python]
class TaskCategory(str, Enum):
    FORECASTING = "forecasting"
    CLASSIFICATION = "classification"
    REGRESSION = "regression"
    CLUSTERING = "clustering"
    ANOMALY_DETECTION = "anomaly_detection"
    DESCRIPTIVE = "descriptive"
    OTHER = "other"
\end{lstlisting}

\subsection{StageStatus}
\begin{lstlisting}[language=Python]
class StageStatus(str, Enum):
    PENDING = "pending"
    RUNNING = "running"
    COMPLETED = "completed"
    FAILED = "failed"
    SKIPPED = "skipped"
\end{lstlisting}

\section{Stage 1 Models}
\subsection{ColumnSummary}
Captures metadata for a single column:
\begin{itemize}
    \item \texttt{name}: Column name
    \item \texttt{dtype}: Pandas data type
    \item \texttt{logical\_type}: Inferred logical type (numeric, categorical, datetime, etc.)
    \item \texttt{null\_fraction}: Fraction of null values
    \item \texttt{n\_unique}: Number of unique values
    \item \texttt{min\_value}, \texttt{max\_value}: For numeric columns
    \item \texttt{unique\_values}: All values if categorical with $\leq$ 20 unique values
\end{itemize}

\subsection{DatasetSummary}
Complete summary of a dataset:
\begin{itemize}
    \item \texttt{filename}, \texttt{filepath}: File identification
    \item \texttt{n\_rows}, \texttt{n\_cols}: Dataset dimensions
    \item \texttt{columns}: List of \texttt{ColumnSummary} objects
    \item \texttt{candidate\_keys}: Potential primary key columns
    \item \texttt{data\_quality\_score}: Overall quality assessment
\end{itemize}

\section{Stage 2 Models}
\subsection{TaskProposal}
A proposed analytical task:
\begin{itemize}
    \item \texttt{id}: Unique identifier (e.g., TSK-001)
    \item \texttt{category}: Task type (forecasting, classification, etc.)
    \item \texttt{title}: Short descriptive title
    \item \texttt{problem\_statement}: Detailed description
    \item \texttt{required\_datasets}: List of required data files
    \item \texttt{target\_column}: Column to predict
    \item \texttt{feature\_columns}: Columns to use as features
    \item \texttt{join\_plan}: For multi-dataset tasks
    \item \texttt{feasibility\_score}: Confidence in task feasibility
    \item \texttt{forecast\_horizon}: For forecasting tasks
\end{itemize}

\section{Stage 3 Models}
\subsection{Stage3Plan}
Detailed execution plan:
\begin{itemize}
    \item \texttt{plan\_id}: Plan identifier (e.g., PLAN-TSK-001)
    \item \texttt{file\_instructions}: How to load each data file
    \item \texttt{join\_steps}: Data joining operations
    \item \texttt{feature\_engineering}: Feature creation specifications
    \item \texttt{target\_column}: Column to predict
    \item \texttt{validation\_strategy}: Train/test splitting strategy
    \item \texttt{evaluation\_metrics}: Metrics to compute
\end{itemize}

\section{Stage 3.5 Models}
\subsection{ForecastingMethod}
A proposed forecasting/ML method:
\begin{itemize}
    \item \texttt{method\_id}: Unique identifier (M1, M2, M3)
    \item \texttt{name}: Method name
    \item \texttt{category}: baseline, statistical, or ml
    \item \texttt{implementation\_code}: Full Python implementation
    \item \texttt{hyperparameters}: Model parameters
\end{itemize}

\subsection{MethodBenchmarkResult}
Benchmark results for a method:
\begin{itemize}
    \item \texttt{method\_id}: Method identifier
    \item \texttt{iterations}: List of benchmark runs
    \item \texttt{average\_metrics}: Mean performance metrics
    \item \texttt{is\_valid}: Whether results are trustworthy
\end{itemize}

\section{Pipeline State}
\subsection{PipelineState}
Overall pipeline state tracking:
\begin{itemize}
    \item \texttt{session\_id}: Unique session identifier
    \item \texttt{selected\_task\_id}: Current task being executed
    \item \texttt{stages}: Dictionary of stage states
    \item \texttt{stage1\_output} through \texttt{stage5\_output}: Stage results
    \item Methods for marking stages started/completed/failed
\end{itemize}

%==============================================================================
% CHAPTER 5: STAGE AGENTS
%==============================================================================
\chapter{Pipeline Stage Agents}

\section{Agent Architecture}
Each stage agent follows a consistent architecture using LangGraph:

\begin{lstlisting}[language=Python]
class StageXState(TypedDict):
    messages: Annotated[list, add_messages]
    # Stage-specific state fields
    iteration: int
    complete: bool

def create_stageX_agent():
    # Create LLM with tools
    llm = ChatOpenAI(**LLM_CONFIG).bind_tools(STAGEX_TOOLS)
    
    def agent_node(state):
        # Agent reasoning logic
        response = llm.invoke(state["messages"])
        return {"messages": [response]}
    
    def should_continue(state):
        # Determine next action or end
        ...
    
    # Build graph
    graph = StateGraph(StageXState)
    graph.add_node("agent", agent_node)
    graph.add_node("tools", ToolNode(STAGEX_TOOLS))
    # Add edges and compile
    return graph.compile()
\end{lstlisting}

\section{Stage 1: Dataset Summarization}
\subsection{Purpose}
Profiles all available datasets and generates comprehensive summaries describing data characteristics, quality, and analysis potential.

\subsection{Tools}
\begin{itemize}
    \item \texttt{list\_available\_datasets()}: List all data files
    \item \texttt{profile\_dataset(filename)}: Generate detailed column-level summary
    \item \texttt{save\_dataset\_summary(filename, summary\_json)}: Persist summary
    \item \texttt{get\_existing\_summaries()}: List already-profiled datasets
    \item \texttt{analyze\_dataset\_for\_forecasting(filename)}: Assess forecasting suitability
\end{itemize}

\subsection{Output}
\texttt{Stage1Output} containing list of \texttt{DatasetSummary} objects saved to \texttt{summaries/} directory.

\section{Stage 2: Task Proposal Generation}
\subsection{Purpose}
Explores dataset summaries and proposes 3-5 analytical tasks, focusing on forecasting when possible.

\subsection{Tools}
\begin{itemize}
    \item \texttt{list\_dataset\_summaries()}: Get available summaries
    \item \texttt{read\_dataset\_summary(filename)}: Read specific summary
    \item \texttt{explore\_data\_relationships()}: Find potential dataset joins
    \item \texttt{evaluate\_forecasting\_feasibility()}: Assess task feasibility
    \item \texttt{save\_task\_proposals(proposals\_json)}: Save proposals
    \item \texttt{get\_proposal\_template()}: Get JSON template
    \item \texttt{analyze\_column\_deeply()}: Deep column analysis
    \item \texttt{find\_best\_target\_columns()}: Rank potential targets
\end{itemize}

\subsection{Output}
\texttt{Stage2Output} with list of \texttt{TaskProposal} objects saved to \texttt{stage2\_out/task\_proposals.json}.

\section{Stage 3: Execution Planning}
\subsection{Purpose}
Creates detailed execution plans for selected tasks, specifying data loading, transformation, and processing steps.

\subsection{Tools}
\begin{itemize}
    \item \texttt{load\_task\_proposal(task\_id)}: Load specific task
    \item \texttt{list\_all\_proposals()}: List available tasks
    \item \texttt{inspect\_data\_file\_stage3(filename)}: Examine data structure
    \item \texttt{validate\_columns\_for\_task()}: Check column data quality
    \item \texttt{analyze\_join\_feasibility()}: Assess multi-dataset joins
    \item \texttt{save\_stage3\_plan(plan\_json)}: Save execution plan
    \item \texttt{validate\_target\_for\_modeling()}: Verify target suitability
\end{itemize}

\subsection{Output}
\texttt{Stage3Plan} saved to \texttt{stage3\_out/PLAN-TSK-XXX.json}.

\section{Stage 3B: Data Preparation}
\subsection{Purpose}
Prepares data according to the execution plan, handling joins, feature engineering, and missing value treatment.

\subsection{Key Operations}
\begin{itemize}
    \item Load datasets according to file instructions
    \item Execute join operations
    \item Apply feature engineering transformations
    \item Handle missing values (imputation/removal)
    \item Create train/validation/test splits
    \item Save prepared data as parquet files
\end{itemize}

\subsection{Output}
\texttt{PreparedDataOutput} with path to prepared parquet file and data quality report.

\section{Stage 3.5A: Method Proposal}
\subsection{Purpose}
Proposes exactly 3 forecasting/ML methods to benchmark, ranging from simple baselines to advanced approaches.

\subsection{Method Categories}
\begin{itemize}
    \item \textbf{Baseline}: Simple methods (mean, last value, naive)
    \item \textbf{Statistical}: ARIMA, exponential smoothing
    \item \textbf{Machine Learning}: Random Forest, XGBoost, LightGBM
\end{itemize}

\subsection{Output}
\texttt{MethodProposalOutput} with 3 \texttt{ForecastingMethod} objects including full implementation code.

\section{Stage 3.5B: Method Benchmarking}
\subsection{Purpose}
Benchmarks proposed methods with multiple iterations to ensure consistency and select the best performer.

\subsection{Process}
\begin{enumerate}
    \item Load method proposals and prepared data
    \item Execute each method 3 times
    \item Calculate performance metrics (MAE, RMSE, MAPE, R²)
    \item Compute coefficient of variation for stability
    \item Select best method based on primary metric
    \item Save model checkpoints
\end{enumerate}

\subsection{Output}
\texttt{TesterOutput} with benchmark results and selected method information.

\section{Stage 4: Execution}
\subsection{Purpose}
Executes the selected method on full data and generates final predictions.

\subsection{Process}
\begin{enumerate}
    \item Load winning method code from Stage 3.5B
    \item Load prepared data
    \item Train final model
    \item Generate predictions
    \item Calculate final metrics
    \item Save predictions and model artifacts
\end{enumerate}

\subsection{Output}
\texttt{ExecutionResult} with predictions file path and performance metrics.

\section{Stage 5: Visualization}
\subsection{Purpose}
Creates insightful visualizations from results and generates interpretation.

\subsection{Visualization Types}
\begin{itemize}
    \item Actual vs Predicted plots
    \item Residual analysis
    \item Feature importance (if applicable)
    \item Time series decomposition
    \item Distribution comparisons
    \item Correlation heatmaps
\end{itemize}

\subsection{Output}
\texttt{VisualizationReport} with list of generated visualizations and insights.

\section{Stage 6: Report Generation}
\subsection{Purpose}
Generates comprehensive final reports synthesizing all pipeline outputs.

\subsection{Report Sections}
\begin{itemize}
    \item Executive Summary
    \item Data Description
    \item Methodology
    \item Results and Metrics
    \item Visualizations
    \item Conclusions and Recommendations
\end{itemize}

\subsection{Output}
JSON report saved to \texttt{stage6\_out/} with embedded visualization paths.

\section{Stage 7: Guardrails Validation}
\subsection{Purpose}
Optional validation stage that runs statistical tests to verify model predictions.

\subsection{Validation Tests}
\begin{itemize}
    \item Correlation analysis between predictions and actuals
    \item Propensity score analysis
    \item Residual normality tests
    \item Heteroscedasticity checks
    \item Out-of-distribution detection
\end{itemize}

\subsection{Output}
Validation report with pass/fail status and detailed test results.

%==============================================================================
% CHAPTER 6: SPECIAL AGENTS
%==============================================================================
\chapter{Special Agents}

\section{Conversation Agent}
\subsection{Purpose}
Handles user queries, determines intent, and orchestrates pipeline stages.

\subsection{Capabilities}
\begin{itemize}
    \item Natural language understanding
    \item Task ID extraction from user input
    \item Pipeline action detection
    \item Session state management
    \item Context-aware responses
\end{itemize}

\subsection{State}
\begin{lstlisting}[language=Python]
class ConversationState(TypedDict):
    messages: Annotated[list, add_messages]
    session_id: str
    user_intent: str
    current_task_id: Optional[str]
    pipeline_action: Optional[str]
    pipeline_stages: List[str]
    response_ready: bool
    iteration: int
\end{lstlisting}

\section{EDA Agent}
\subsection{Purpose}
Handles exploratory data analysis queries by writing and executing custom Python code.

\subsection{Capabilities}
\begin{itemize}
    \item Understanding natural language data queries
    \item Writing custom analysis code
    \item Creating visualizations on demand
    \item Detecting new datasets
    \item Generating insights from data
\end{itemize}

\subsection{Tools}
\begin{itemize}
    \item \texttt{write\_and\_execute\_code()}: Execute analysis code
    \item \texttt{create\_visualization()}: Generate plots
    \item \texttt{scan\_for\_new\_datasets()}: Find new data files
    \item \texttt{get\_dataset\_info()}: Quick dataset overview
    \item \texttt{save\_eda\_report()}: Persist analysis results
\end{itemize}

%==============================================================================
% CHAPTER 7: ORCHESTRATION
%==============================================================================
\chapter{Master Orchestrator}

\section{Overview}
The \texttt{master\_orchestrator.py} module coordinates pipeline execution, manages state transitions, and handles the conversational interface.

\section{Pipeline Builder}
\begin{lstlisting}[language=Python]
def build_pipeline_graph(start_stage="stage1", end_stage="stage6"):
    """Build a pipeline graph from start to end stage."""
    graph = StateGraph(PipelineState)
    
    # Add stage nodes
    stages_to_run = STAGE_ORDER[
        STAGE_ORDER.index(start_stage):
        STAGE_ORDER.index(end_stage) + 1
    ]
    
    for stage in stages_to_run:
        graph.add_node(stage, STAGE_NODES[stage])
    
    # Add sequential edges
    for i, stage in enumerate(stages_to_run[:-1]):
        graph.add_edge(stage, stages_to_run[i+1])
    
    return graph.compile()
\end{lstlisting}

\section{State Loading}
The orchestrator supports checkpoint-based resumption:
\begin{itemize}
    \item Loads cached state from disk
    \item Determines where to resume execution
    \item Skips completed stages
    \item Handles partial failures
\end{itemize}

\section{ConversationalOrchestrator}
\begin{lstlisting}[language=Python]
class ConversationalOrchestrator:
    """Orchestrates conversational pipeline experience."""
    
    def __init__(self, session_id=None):
        self.session_id = session_id or generate_session_id()
        self.handler = ConversationHandler(self.session_id)
    
    def process_user_input(self, user_input: str):
        """Process user input and return response."""
        # Check for EDA queries
        if self._is_eda_query(user_input):
            return self._handle_eda_query(user_input)
        
        # Process through conversation handler
        result = self.handler.process_message(user_input)
        
        # Execute pipeline if action detected
        if result.get("action") == "run_pipeline":
            self._execute_pipeline(result["task_id"])
        
        return result
\end{lstlisting}

%==============================================================================
% CHAPTER 8: WEB SERVER
%==============================================================================
\chapter{Web Server}

\section{Overview}
The \texttt{server.py} module provides a FastAPI-based web interface with WebSocket support for real-time updates.

\section{API Endpoints}
\begin{longtable}{|l|l|p{6cm}|}
\hline
\textbf{Method} & \textbf{Endpoint} & \textbf{Description} \\
\hline
POST & /chat & Process chat message \\
\hline
GET & /status & Get pipeline status \\
\hline
GET & /stages & Get all stage statuses \\
\hline
GET & /stages/\{stage\} & Get specific stage output \\
\hline
GET & /stages/\{stage\}/files & List stage output files \\
\hline
GET & /proposals & Get task proposals \\
\hline
POST & /run/\{task\_id\} & Execute pipeline for task \\
\hline
GET & /conversations & List conversation sessions \\
\hline
GET & /conversations/\{id\} & Get conversation history \\
\hline
POST & /guardrails/\{plan\_id\} & Run guardrails validation \\
\hline
WS & /ws/task-progress & Real-time progress updates \\
\hline
\end{longtable}

\section{Pipeline Tracker}
\begin{lstlisting}[language=Python]
class PipelineTracker:
    """Tracks pipeline execution state for UI."""
    
    def __init__(self):
        self.current_task_id = None
        self.is_running = False
        self.current_stage = None
        self.stage_status = {}
        self.errors = []
        self.lock = threading.Lock()
    
    def start_pipeline(self, task_id):
        with self.lock:
            self.current_task_id = task_id
            self.is_running = True
    
    def update_stage(self, stage, status):
        with self.lock:
            self.current_stage = stage
            self.stage_status[stage] = status
\end{lstlisting}

\section{Startup Prerequisites}
The server ensures prerequisites before accepting requests:
\begin{enumerate}
    \item All datasets have summaries (runs Stage 1 if needed)
    \item 3-5 task proposals exist (runs Stage 2 if needed)
\end{enumerate}

%==============================================================================
% CHAPTER 9: UTILITIES
%==============================================================================
\chapter{Utility Functions}

\section{Data Profiling}
\begin{itemize}
    \item \texttt{infer\_logical\_type(series)}: Infer column type from pandas Series
    \item \texttt{get\_column\_statistics(series, type)}: Get type-appropriate statistics
    \item \texttt{analyze\_column\_semantics()}: Understand column meaning
    \item \texttt{profile\_csv(filepath)}: Generate complete dataset profile
\end{itemize}

\section{File Operations}
\begin{itemize}
    \item \texttt{list\_data\_files(directory)}: List available data files
    \item \texttt{load\_dataframe(filepath)}: Load DataFrame from various formats
    \item \texttt{inspect\_data\_file(filename)}: Get formatted file inspection
\end{itemize}

\section{Python Sandbox}
The \texttt{PythonSandbox} class provides safe code execution:
\begin{lstlisting}[language=Python]
class PythonSandbox:
    """Safe Python code execution environment."""
    
    DEFAULT_IMPORTS = {
        'pd': pandas,
        'np': numpy,
        'plt': matplotlib.pyplot,
        'json': json,
    }
    
    @classmethod
    def execute(cls, code, additional_globals=None):
        """Execute code with output capture."""
        namespace = cls.DEFAULT_IMPORTS.copy()
        if additional_globals:
            namespace.update(additional_globals)
        
        # Capture stdout
        output = io.StringIO()
        with redirect_stdout(output):
            exec(code, namespace)
        
        return True, output.getvalue(), namespace
\end{lstlisting}

\section{JSON Parsing}
Robust JSON parsing with multiple fallback strategies:
\begin{itemize}
    \item \texttt{extract\_json\_block(text)}: Extract JSON from mixed text
    \item \texttt{parse\_tool\_call(raw)}: Parse LLM tool calls
    \item \texttt{parse\_proposals\_json(raw)}: Parse task proposals
\end{itemize}

%==============================================================================
% CHAPTER 10: USER INTERFACE
%==============================================================================
\chapter{User Interface}

\section{HTML Pages}
\begin{itemize}
    \item \texttt{index.html}: Main chat interface
    \item \texttt{outputs.html}: View stage outputs
    \item \texttt{logs.html}: Pipeline execution logs
    \item \texttt{guardrails.html}: Guardrails validation interface
    \item \texttt{task\_details.html}: Detailed task view
\end{itemize}

\section{JavaScript Modules}
\begin{itemize}
    \item \texttt{chat.js}: Chat functionality
    \item \texttt{api.js}: API communication
    \item \texttt{websocket.js}: Real-time updates
    \item \texttt{outputs.js}: Output display
    \item \texttt{stages.js}: Stage visualization
\end{itemize}

%==============================================================================
% CHAPTER 11: RUNNING THE SYSTEM
%==============================================================================
\chapter{Running the System}

\section{Command Line Interface}
\begin{lstlisting}[language=bash]
# Interactive conversation mode
python run_conversational.py

# Quick analysis mode
python run_conversational.py --analyze

# Propose tasks for a query
python run_conversational.py --propose "forecast sales"

# Execute specific task
python run_conversational.py --execute TSK-001

# Run full pipeline
python run_conversational.py --full TSK-001

# Show configuration
python run_conversational.py --config
\end{lstlisting}

\section{Web Server}
\begin{lstlisting}[language=bash]
# Start the web server
python server.py

# Access at http://localhost:8000
\end{lstlisting}

\section{Environment Variables}
\begin{itemize}
    \item \texttt{LLM\_BASE\_URL}: LLM API endpoint
    \item \texttt{LLM\_API\_KEY}: API authentication key
    \item \texttt{USE\_GROQ}: Whether to use Groq API
    \item \texttt{GROQ\_API\_KEY}: Groq API key if using Groq
    \item \texttt{LANGSMITH\_TRACING}: Enable LangSmith tracing
\end{itemize}

%==============================================================================
% CHAPTER 12: DATA FLOW
%==============================================================================
\chapter{Data Flow}

\section{Stage-to-Stage Data Passing}
\begin{enumerate}
    \item \textbf{Stage 1 $\rightarrow$ Stage 2}: Dataset summaries (JSON)
    \item \textbf{Stage 2 $\rightarrow$ Stage 3}: Task proposals (JSON)
    \item \textbf{Stage 3 $\rightarrow$ Stage 3B}: Execution plan (JSON)
    \item \textbf{Stage 3B $\rightarrow$ Stage 3.5A}: Prepared data (Parquet)
    \item \textbf{Stage 3.5A $\rightarrow$ Stage 3.5B}: Method proposals (JSON)
    \item \textbf{Stage 3.5B $\rightarrow$ Stage 4}: Selected method + checkpoints
    \item \textbf{Stage 4 $\rightarrow$ Stage 5}: Predictions (Parquet/CSV)
    \item \textbf{Stage 5 $\rightarrow$ Stage 6}: Visualizations + insights
    \item \textbf{Stage 6 $\rightarrow$ Stage 7}: Final report for validation
\end{enumerate}

\section{Artifact Storage}
All artifacts are stored with:
\begin{itemize}
    \item Unique artifact IDs
    \item SHA256 checksums for integrity
    \item Metadata sidecar files
    \item Atomic write operations
\end{itemize}

%==============================================================================
% APPENDIX
%==============================================================================
\appendix
\chapter{Tool Reference}

\section{Stage 1 Tools}
\begin{itemize}
    \item \texttt{list\_available\_datasets()}
    \item \texttt{profile\_dataset(filename)}
    \item \texttt{save\_dataset\_summary(filename, summary\_json)}
    \item \texttt{get\_existing\_summaries()}
    \item \texttt{analyze\_dataset\_for\_forecasting(filename)}
\end{itemize}

\section{Stage 2 Tools}
\begin{itemize}
    \item \texttt{list\_dataset\_summaries()}
    \item \texttt{read\_dataset\_summary(filename)}
    \item \texttt{explore\_data\_relationships(summary\_files)}
    \item \texttt{python\_sandbox\_stage2(code, description)}
    \item \texttt{evaluate\_forecasting\_feasibility(target, date, summaries)}
    \item \texttt{save\_task\_proposals(proposals\_json)}
    \item \texttt{get\_proposal\_template()}
    \item \texttt{analyze\_column\_deeply(dataset, column)}
    \item \texttt{find\_best\_target\_columns(dataset)}
\end{itemize}

\section{Stage 3 Tools}
\begin{itemize}
    \item \texttt{load\_task\_proposal(task\_id)}
    \item \texttt{list\_all\_proposals()}
    \item \texttt{list\_data\_files\_stage3()}
    \item \texttt{inspect\_data\_file\_stage3(filename, n\_rows)}
    \item \texttt{validate\_columns\_for\_task(filename, columns)}
    \item \texttt{analyze\_join\_feasibility(left, right, left\_key, right\_key)}
    \item \texttt{python\_sandbox\_stage3(code, description)}
    \item \texttt{save\_stage3\_plan(plan\_json)}
    \item \texttt{get\_execution\_plan\_template()}
\end{itemize}

\end{document}
